\documentclass[11pt,letterpaper]{article}
\usepackage{physical_ai_legislation}
\usepackage[backend=biber,style=numeric,sorting=none]{biblatex}
\addbibresource{physical_ai_emergency.bib}

\title{%
  \textbf{Physical AI Oncology Clinical Trial Site}\\
  \textbf{Emergency Preparedness Plan}\\[6pt]
  \large Emergency Response, Disaster Recovery,\\
  and Business Continuity for Physical AI Operations\\[6pt]
  \normalsize Comprehensive emergency preparedness covering robot system\\
  failures, cybersecurity incidents, natural disasters, and clinical emergencies.
}
\author{%
  Kevin Kawchak\\
  CEO, ChemicalQDevice\\
  \texttt{kevink@chemicalqdevice.com}
}
\date{March 24, 2026}

\begin{document}
\maketitle
\thispagestyle{fancy}

\begin{center}
\footnotesize
\textit{This draft is independent and is not endorsed, sponsored, or approved
by any trial sponsor, CRO, site, IRB, regulator, or medical society;
and was adapted using Claude Code Opus 4.6.}
\end{center}

\vspace{1em}
\tableofcontents
\newpage

\section{Emergency Classification System}

Physical AI oncology clinical trial sites shall classify emergencies into four
levels, each with defined response protocols:

\subsection{Level 1 -- Minor Incident}

\begin{enumerate}[label=(\alph*)]
\item Definition: An event affecting a single Physical AI system or a single
  patient without immediate risk of serious harm. Examples include a single
  robot entering fault mode, a Grade 1 adverse event, or a minor network
  disruption affecting non-critical systems.

\item Response: The affected Physical AI system is automatically placed in safe
  mode. The site safety officer assesses the situation and determines whether
  the system can be restarted or requires maintenance. Patient care continues
  using alternative systems. The simulation demonstrated that the fleet of 29
  robot instances provides redundancy within each system category, allowing
  continued operations when individual systems are
  unavailable~\cite{newtrial2026}.

\item Reporting: Documented in the shift log and included in the monthly report
  to CDPH.
\end{enumerate}

\subsection{Level 2 -- Moderate Emergency}

\begin{enumerate}[label=(\alph*)]
\item Definition: An event affecting multiple Physical AI systems, multiple
  patients, or requiring activation of emergency medical response. Examples
  include multiple simultaneous robot faults, a Grade 2 or Grade 3 adverse
  event, a localized power failure, or a cybersecurity anomaly.

\item Response: The site safety officer activates the emergency response team.
  The on-call emergency physician is notified. Affected clinical areas are
  secured. Patient procedures in unaffected areas may continue. All active
  procedures in affected areas are safely concluded or aborted using the
  emergency stop protocols described in the patient robot
  instructions~\cite{patientinstructions2026}.

\item Reporting: Reported to CDPH within 72 hours. If a cybersecurity incident
  with potential patient harm, reported within 7 days per the adapted 21 CFR
  Part 312 \S~312.32~\cite{cfr312adapt}.
\end{enumerate}

\subsection{Level 3 -- Major Emergency}

\begin{enumerate}[label=(\alph*)]
\item Definition: An event requiring partial or complete facility evacuation,
  affecting the majority of Physical AI systems, or involving a serious
  adverse event (Grade 4 or higher). Examples include building-wide power
  failure exceeding UPS capacity, structural damage, confirmed cybersecurity
  breach affecting patient safety, or a mass casualty event.

\item Response: The site administrator activates the facility emergency plan.
  All Physical AI systems execute automated safe-shutdown procedures. Patient
  evacuation begins per the evacuation plan. External emergency services are
  activated. The Physical AI Safety Review Committee is convened within 24
  hours~\cite{cfr312adapt}.

\item Reporting: Reported to CDPH and FDA immediately. Fatal or
  life-threatening events reported within 7 days per \S~312.32. Serious
  events reported within 15 days~\cite{cfr312adapt}.
\end{enumerate}

\subsection{Level 4 -- Catastrophic Event}

\begin{enumerate}[label=(\alph*)]
\item Definition: An event rendering the facility completely non-operational,
  such as a major earthquake, widespread infrastructure failure, or complete
  cybersecurity compromise.

\item Response: Complete facility evacuation. All Physical AI systems
  de-energized. Business continuity plan activated. Patient care transferred
  to partner facilities. IND sponsor notified for potential clinical
  hold assessment~\cite{cfr312adapt}.

\item Reporting: Immediate notification to CDPH, FDA, and local emergency
  management. Full incident report within 30 days.
\end{enumerate}

\section{Physical AI System Emergency Procedures}

\subsection{Emergency Stop Protocol}

\begin{enumerate}[label=(\alph*)]
\item Every Physical AI system shall have a clearly marked, accessible emergency
  stop mechanism that can be activated by site safety officers, clinical
  personnel, or patients where clinically appropriate.

\item Upon emergency stop activation, the system shall:
  \begin{enumerate}[label=(\arabic*)]
  \item Immediately cease all autonomous motion.
  \item Retract any instruments or probes that are in contact with or
    inside the patient, using a controlled withdrawal sequence.
  \item Lower treatment couches to the lowest safe position.
  \item De-energize radiation sources (RT systems).
  \item Maintain essential monitoring and life safety functions.
  \item Generate an emergency stop event record in the audit trail with
    timestamp, trigger source, and system state at the time of
    activation~\cite{iche6r3adapt}.
  \end{enumerate}

\item The patient robot instructions for each system type include specific
  emergency behavior; for example, surgical robots retract arms through chest
  ports, RT positioning systems lower couches to exit height, and
  exoskeletons return to standing rest position~\cite{patientinstructions2026}.
\end{enumerate}

\subsection{Cybersecurity Incident Response}

\begin{enumerate}[label=(\alph*)]
\item The cybersecurity incident response plan shall address Physical
  AI-specific threats:
  \begin{enumerate}[label=(\arabic*)]
  \item Robot control system compromise (unauthorized commands to Physical AI
    systems).
  \item AI model poisoning (manipulation of AI model parameters or training
    data).
  \item Digital twin data manipulation (alteration of patient-specific models).
  \item MCP-PAI server infrastructure attack (compromise of authorization,
    FHIR, DICOM, ledger, or provenance servers).
  \item Network segmentation breach (traffic between isolated network
    zones).
  \end{enumerate}

\item Upon detection of a cybersecurity incident:
  \begin{enumerate}[label=(\arabic*)]
  \item All affected Physical AI systems are immediately isolated from the
    network.
  \item Active procedures on affected systems are safely concluded or aborted.
  \item The cybersecurity operations specialist assesses the scope and
    severity.
  \item Patient data integrity is verified using hash-chained audit trail
    validation~\cite{iche6r3adapt}.
  \item The incident is classified and reported per the adapted 21 CFR
    Part 312 timelines~\cite{cfr312adapt}.
  \end{enumerate}

\item The site shall conduct annual penetration testing and quarterly
  vulnerability assessments of all Physical AI system interfaces and
  MCP-PAI infrastructure.
\end{enumerate}

\section{Natural Disaster Preparedness}

\subsection{Seismic Preparedness}

\begin{enumerate}[label=(\alph*)]
\item The facility shall be designed to Seismic Design Category D or higher
  per ASCE 7 and the California Building Code, appropriate for San Francisco's
  seismic zone.

\item Physical AI systems shall be seismically anchored to their mounting
  surfaces with anchorage rated for the systems' operating weight plus 50
  percent margin.

\item Automated seismic sensors shall trigger Physical AI system safe-mode
  procedures when ground acceleration exceeds 0.1g, including retracting
  surgical instruments, lowering couches, and pausing radiation delivery.

\item Post-earthquake inspection protocols shall verify Physical AI system
  calibration before resuming patient procedures.
\end{enumerate}

\subsection{Power Failure}

\begin{enumerate}[label=(\alph*)]
\item The dual utility feed design provides primary power redundancy. If both
  feeds fail, the automatic transfer switch engages emergency generators
  within 10 seconds.

\item UPS systems bridge the transfer gap, providing 30 minutes for surgical
  suites (sufficient to complete or safely abort any procedure) and 15 minutes
  for server rooms (sufficient for orderly AI system shutdown and data
  preservation).

\item If generator fuel is projected to be exhausted before utility power
  restoration, the site shall begin controlled patient discharge and Physical
  AI system shutdown, prioritizing completion of active procedures.

\item The 72-hour generator fuel capacity provides sufficient time for
  coordinated patient transfer to partner facilities in extended outage
  scenarios.
\end{enumerate}

\section{Clinical Emergency Procedures}

\subsection{Medical Emergency During Physical AI Procedure}

\begin{enumerate}[label=(\alph*)]
\item If a patient experiences a medical emergency during a Physical AI-assisted
  procedure:
  \begin{enumerate}[label=(\arabic*)]
  \item The Physical AI system's vital sign monitoring detects the emergency
    and initiates automated response protocols.
  \item Emergency stop is activated automatically or by the site safety officer.
  \item The on-call emergency physician is notified immediately.
  \item The Physical AI system retracts instruments and positions the patient
    for emergency access.
  \item Resuscitation equipment is deployed from the emergency station located
    within each clinical wing.
  \item The event is documented in real time through the Physical AI system's
    audit trail and the patient's electronic health record.
  \end{enumerate}

\item The simulation demonstrated that real-time vital sign monitoring at
  minute-level resolution enables early detection of adverse events; all
  seven adverse events in the simulation were detected, graded, and managed
  through the automated monitoring
  system~\cite{newtrial2026}.
\end{enumerate}

\subsection{Multiple Simultaneous Adverse Events}

\begin{enumerate}[label=(\alph*)]
\item If multiple adverse events occur simultaneously (which did not occur in
  the simulation but is planned for), the site shall:
  \begin{enumerate}[label=(\arabic*)]
  \item Triage events by severity grade, prioritizing the highest-grade event.
  \item Activate the staffing contingency plan to bring additional personnel
    on-site within 30 minutes.
  \item Suspend new patient intake until the simultaneous events are resolved.
  \item Notify the Physical AI Safety Review Committee if the events suggest
    a systemic issue affecting multiple Physical AI systems or patients.
  \end{enumerate}
\end{enumerate}

\section{Business Continuity}

\begin{enumerate}[label=(\alph*)]
\item The site shall maintain mutual aid agreements with at least two hospital
  or clinical facilities within 30 minutes travel time for patient transfer
  during extended emergencies.

\item Critical Physical AI system configuration data, patient records, and audit
  trails shall be backed up to a geographically remote data center with
  recovery point objective (RPO) of 15 minutes and recovery time objective
  (RTO) of 4 hours.

\item The business continuity plan shall include provisions for resuming
  Physical AI operations following a major emergency, including system
  recalibration, Phase 0 re-validation if systems were physically damaged,
  and regulatory notification of operational restart.

\item Annual business continuity exercises shall test the full emergency
  response chain from emergency detection through patient transfer, Physical
  AI system shutdown, data preservation, and operational restart.
\end{enumerate}

\section{Training and Drills}

\begin{enumerate}[label=(\alph*)]
\item All site personnel shall complete emergency preparedness training as part
  of the 160-hour site safety officer certification program.

\item Emergency drills shall be conducted quarterly, rotating through the
  following scenarios:
  \begin{enumerate}[label=(\arabic*)]
  \item Quarter 1: Physical AI system failure during active procedure
    (emergency stop, patient safety, system recovery).
  \item Quarter 2: Cybersecurity incident (network isolation, data integrity
    verification, incident reporting).
  \item Quarter 3: Building evacuation (Physical AI safe-shutdown, patient
    evacuation, external emergency services coordination).
  \item Quarter 4: Extended power failure (generator operation, controlled
    patient discharge, business continuity activation).
  \end{enumerate}

\item Drill performance shall be documented and reviewed by the Physical AI
  Safety Review Committee, with corrective actions for any deficiencies
  identified.

\item The patient journey demonstrated that 10 distinct stages of care, from
  pre-screening through data lock, require coordinated emergency response
  capabilities at each stage~\cite{patientjourney2026}.
\end{enumerate}

\printbibliography[title={References}]

\end{document}
